% pyt-layout.tex - Seitenlayout, Fonts und Farben für die Dokumentation
% ConTeXt MkIV / LuaMetaTeX

% =============================================================================
% SPRACHE (Deutsch)
% =============================================================================

\mainlanguage[de]
\setuplabeltext[de][content=Inhaltsverzeichnis]
\setuplabeltext[de][figure=Abbildung]
\setuplabeltext[de][table=Tabelle]
\setuplabeltext[de][chapter=Kapitel]

% =============================================================================
% SEITENLAYOUT
% =============================================================================

\setuppapersize[A4]

\setuplayout[
    width=middle,
    height=middle,
    topspace=2.5cm,
    backspace=2.5cm,
    header=1cm,
    footer=1cm,
    headerdistance=0.5cm,
    footerdistance=0.5cm,
]

% =============================================================================
% SCHRIFTEN
% =============================================================================

% Moderne serifenlose Schrift für Überschriften
\definefontfamily[mainface][sans][TeX Gyre Heros]
\definefontfamily[mainface][serif][TeX Gyre Pagella]
\definefontfamily[mainface][mono][TeX Gyre Cursor][scale=0.85]

\setupbodyfont[mainface, 11pt]

% =============================================================================
% FARBEN
% =============================================================================

\definecolor[primarycolor][h=2563EB]      % Blau
\definecolor[secondarycolor][h=059669]    % Grün
\definecolor[accentcolor][h=D97706]       % Orange
\definecolor[warningcolor][h=DC2626]      % Rot
\definecolor[codebackground][h=F3F4F6]    % Hellgrau
\definecolor[codeborder][h=D1D5DB]        % Grau

% =============================================================================
% ÜBERSCHRIFTEN
% =============================================================================

\setuphead[chapter][
    style=\bfc\bf,
    color=primarycolor,
    before={\blank[3*big]},
    after={\blank[2*big]},
]

\setuphead[section][
    style=\bfb\bf,
    color=primarycolor,
    before={\blank[2*big]},
    after={\blank[big]},
]

\setuphead[subsection][
    style=\bfa\bf,
    color=primarycolor,
    before={\blank[big]},
    after={\blank[medium]},
]

\setuphead[subsubsection][
    style=\bf,
    before={\blank[medium]},
    after={\blank[small]},
]

% =============================================================================
% GIT-INFORMATIONEN
% =============================================================================

% Git-Branch und Commit-ID dynamisch ermitteln
\startluacode
    local function get_git_info()
        local branch = "unknown"
        local commit = "000000"

        -- Branch ermitteln
        local handle = io.popen("git rev-parse --abbrev-ref HEAD 2>/dev/null")
        if handle then
            branch = handle:read("*a"):gsub("%s+", "")
            handle:close()
        end

        -- Kurze Commit-ID ermitteln
        handle = io.popen("git rev-parse --short HEAD 2>/dev/null")
        if handle then
            commit = handle:read("*a"):gsub("%s+", "")
            handle:close()
        end

        return branch, commit
    end

    local branch, commit = get_git_info()
    document.gitbranch = branch
    document.gitcommit = commit
    document.gitinfo = "[git] • Branch: " .. branch .. "@" .. commit
\stopluacode

% Makro für Git-Info in TeX verfügbar machen
\def\gitinfo{\ctxlua{context(document.gitinfo)}}
\def\gitbranch{\ctxlua{context(document.gitbranch)}}
\def\gitcommit{\ctxlua{context(document.gitcommit)}}

% =============================================================================
% KOPF- UND FUSSZEILEN
% =============================================================================

% Markierungen für Kapitel und Section aktivieren
\setuphead[chapter][marking=chapter]
\setuphead[section][marking=section]

% Seitennummerierung NUR in Fusszeile (nicht in Kopfzeile)
\setuppagenumbering[location=]

% Kopfzeile: Links = Dokumenttitel, Rechts = Kapitel (uppercase)
\setupheadertexts
    [Python mit KI][\WORD{\getmarking[chapter]}]

% Fusszeile: Git-Info zentriert, Seitenzahl rechts
% Drei-Spalten-Layout: leer | Git-Info | Seitenzahl
\setupfootertexts
    [][\hbox to \hsize{\hfil\gitinfo\hfil\llap{\pagenumber}}]

\setupheader[style=\small, color=black]
\setupfooter[style=\small, color=black]

% Titelseite ohne Kopf-/Fusszeile
\definepagebreak[titelpagebreak][yes,header,footer]

% =============================================================================
% ABSÄTZE UND ABSTÄNDE
% =============================================================================

\setupwhitespace[medium]
\setupindenting[no]

% =============================================================================
% LISTEN
% =============================================================================

\setupitemize[each][packed][
    margin=1em,
    before={\blank[small]},
    after={\blank[small]},
]

% =============================================================================
% LINKS
% =============================================================================

\setupinteraction[
    state=start,
    color=primarycolor,
    contrastcolor=primarycolor,
    style=normal,
]

% =============================================================================
% INHALTSVERZEICHNIS
% =============================================================================

\setupcombinedlist[content][
    list={chapter,section},
    alternative=c,
]

\setuplist[chapter][
    style=\bf,
    before={\blank[small]},
]

\setuplist[section][
    margin=1.5em,
]

